\scanpage{31}
\mspage{27}
% archive sheet number: 30

\uline{Digression 2}\quad Notion de sous-magasin.

Soit $\unsure{M|\mathcal M}$, \inserted{$\trianglelefteq,\overset{\circ}{\ll},|\circ|$} un magasin, $M'$ une partie, munissons la des relations induites -- quand peut-on dire que c'est un magasin.

Mag 1 trivial

Mag $2'$ a) aussi. Pour b, il faut supposer

[\uuline{SousMag}] Si $X\overset{\circ}{\ll}Y$ dans $M'$, et $X'\unsure{\triangleleft|\trianglelefteq}X$ dans $M'$, alors l'unique $Y'\in\widetilde Y$ tel que $X'\overset{\circ}{\ll}Y'$ est $\in M'$.
% slanted
\margin{\uline{NB} Il faudrait imposer une condition plus forte, à voir. \uline{Sousmag} Si $X\overset{\circ}{\ll}Y'\unsure{\triangleleft|\trianglelefteq}Y$, et $X,Y\in M'$, alors $Y'\in M'$, laquelle s'écrit (pour \inserted{$X,Y\in M'$}) : si $X\unsure{\overset{\circ}{\ll}|\ll}Y$ dans $M$, \deleted{alors} \inserted{alors $X\underset{M'}{\ll}Y$}. Les \unsure{couples plong. ds $M'$ sont alors compatibles|couples plong. ds $M'$ sont alors complets}. \unsure{serait plus forte}.}

(\uline{NB} c'est une figure si $M'$ fermé par $\trianglelefteq$)

Mag $3'$ trivial, Mag $4'$ trivial.

\uline{Ex 1} Soit $\unsure{\Phi|\varphi}$ \inserted{$\subset M$} une préfigure de $M$, i.e. une partie telle que $\overline\Phi^{(\trianglelefteq)}$ soit une figure, ou encore, \deleted{$\Phi$ fermé} $\forall X,Y\in\Phi$, on a $X\unsure{\gtrless|\lessgtr|\bowtie}Y$.

Considérons $\operatorname{Omb}^{0}(\Phi)$ (attention, c'est fermé pour $\overset{\circ}{\ll}$, mais pas en général pour $\trianglelefteq$). \deleted{\unsure{Dans ce cas}} \inserted{Alors} $M'=\operatorname{Omb}^{0}(\Phi)$ est un sous-magasin, i.e. satisfait Sousmag

\deleted{est un sous-magasin, i.e. satisfait la condition Sous-mag ci-dessus, il faut \unsure{s.} que l'on ait}
\[
\text{\deleted{\uline{Sousmag $\Phi$}}}\quad
\deleted{\forall Z\in\Phi\text{ et }X\overset{\circ}{\ll}Z,\text{ soit }X'\in\widetilde X,}
\]
\deleted{et $Z'$ son image dans $\widetilde Z_{\unsure{a|\alpha}}$ (donc $Z'\in\overline\Phi$), on veut que $X'\in\operatorname{Omb}^{0}(\Phi)\Rightarrow Z'\in\operatorname{Omb}^{0}(\Phi)$ i.e. en fait \uline{$Z'\in\Phi$ (puisque $Z'\in\overline\Phi$)} (même s'il n'est pas fermé par $\trianglelefteq$, quand $\Phi$ n'est pas une figure\ldots)}
